\documentclass[handout]{beamer}
\mode<presentation>
\usepackage{xcolor}
\usepackage[toc]{multitoc}
\usepackage{color}
\usepackage{graphicx}
\graphicspath{{../buku_ajar/buku_ajar_logika_matematika/figures/}}
\usepackage{tikz}
\usetikzlibrary{shapes,arrows,positioning,calc}
\usepackage{listings}
\usepackage{lmodern}
\usepackage{amsmath}
\usetheme[secheader]{Boadilla}
\usecolortheme{whale}
\setbeamertemplate{caption}[numbered]
\setbeamertemplate{bibliography item}[text]
\setbeamertemplate{navigation symbols}{}

\newcommand{\logic}[1]{\textbf{\textit{#1}}}

\theoremstyle{definition}
\newtheorem{definisi}[theorem]{Definisi}
\theoremstyle{example}
\newtheorem{contoh}[theorem]{Contoh}

\renewcommand*{\multicolumntoc}{2}

\setbeamertemplate{footline}
{
  \leavevmode%
  \hbox{%
  \begin{beamercolorbox}[wd=.333333\paperwidth,ht=2.25ex,dp=1ex,center]{author in head/foot}%
    \usebeamerfont{author in head/foot}\insertshortauthor%
  \end{beamercolorbox}%
  \begin{beamercolorbox}[wd=.433333\paperwidth,ht=2.25ex,dp=1ex,center]{title in head/foot}%
    \usebeamerfont{title in head/foot}\insertshorttitle
  \end{beamercolorbox}%
  \begin{beamercolorbox}[wd=.233333\paperwidth,ht=2.25ex,dp=1ex,right]{date in head/foot}%
    \usebeamerfont{date in head/foot}\insertshortdate{}\hspace*{2em} \insertframenumber{} / \inserttotalframenumber\hspace*{2ex} 
  \end{beamercolorbox}}%
  \vskip0pt%
}

\subtitle[Logika Matematika]{Logika Matematika}
\title[Aljabar Boolean]{Prinsip dan Hukum-hukum Dasar Aljabar Boolean}
\author{Cecep Suwanda, S.Si., M.Kom.}
\date{Maret 2025}
\institute{Universitas Bale Bandung}

\begin{document}

\maketitle

\section*{Daftar Isi}
\begin{frame}
\frametitle{Daftar Isi}
\tableofcontents
\end{frame}

% ============================================================
% SECTION 9-1: Pendahuluan
% ============================================================
\section{Pendahuluan Aljabar Boolean}

\begin{frame}
\frametitle{Sejarah Aljabar Boolean}
\begin{itemize}
  \item \textbf{George Boole} (1854): \textit{An Investigation of the Laws of Thought}. Aljabar untuk penalaran logis.
  \item \textbf{Claude Shannon} (1938): \textit{A Symbolic Analysis of Relay and Switching Circuits}. Aplikasi ke rangkaian sakelar.
\end{itemize}
\begin{block}{Sistem}
Himpunan $\{0, 1\}$ dengan tiga operasi: OR ($+$), AND ($\cdot$), NOT ($'$). Fondasi komputasi digital.
\end{block}
\end{frame}

\begin{frame}
\frametitle{Relasi Isomorfik}
Logika Proposisi $\leftrightarrow$ Teori Himpunan $\leftrightarrow$ Aljabar Boolean
\begin{center}
\small
\begin{tabular}{|l|l|l|l|}
\hline
\textbf{Konsep} & \textbf{Logika} & \textbf{Himpunan} & \textbf{Boolean} \\ \hline
Benar & T & $U$ & 1 \\ \hline
Salah & F & $\emptyset$ & 0 \\ \hline
AND & $\wedge$ & $\cap$ & $\cdot$ \\ \hline
OR & $\vee$ & $\cup$ & $+$ \\ \hline
NOT & $\neg$ & $^c$ & $'$ \\ \hline
\end{tabular}
\end{center}
De Morgan: $\neg(p \wedge q) \equiv \neg p \vee \neg q$ $\Leftrightarrow$ $(xy)' = x' + y'$
\end{frame}

\begin{frame}
\frametitle{Definisi Aljabar Boolean}
\begin{definisi}[Aljabar Boolean]
Sistem $\langle B, +, \cdot, ', 0, 1 \rangle$: himpunan $B$, operator biner $+$ dan $\cdot$, operator uner $'$, elemen identitas 0 dan 1, yang memenuhi Postulat Huntington.
\end{definisi}
\end{frame}

% ============================================================
% SECTION 9-2: Postulat Huntington
% ============================================================
\section{Postulat Huntington}

\begin{frame}
\frametitle{Enam Postulat Huntington (1904)}
\begin{enumerate}
  \item \textbf{Ketertutupan}: $a+b \in B$, $a \cdot b \in B$
  \item \textbf{Identitas}: $a+0=a$, $a \cdot 1 = a$
  \item \textbf{Komutatif}: $a+b=b+a$, $a \cdot b = b \cdot a$
  \item \textbf{Distributif}: $a \cdot (b+c) = ab + ac$; $a + (bc) = (a+b)(a+c)$
  \item \textbf{Komplemen}: $a + a' = 1$, $a \cdot a' = 0$
  \item \textbf{Eksistensi minimal}: $0 \neq 1$
\end{enumerate}
Hukum asosiatif \textbf{bukan} postulat; merupakan teorema turunan.
\end{frame}

\begin{frame}
\frametitle{Boolean vs Aljabar Konvensional}
\begin{center}
\small
\begin{tabular}{|l|l|l|}
\hline
\textbf{Aspek} & \textbf{Konvensional} & \textbf{Boolean} \\ \hline
Domain & $\mathbb{R}$ & $\{0,1\}$ \\ \hline
Distributif & Hanya $\times$ atas $+$ & Keduanya ($+$ atas $\cdot$ juga) \\ \hline
Idempoten & $a+a \neq a$ & $a+a=a$, $a \cdot a = a$ \\ \hline
Invers & Ada ($-a$, $a^{-1}$) & Tidak ada; hanya komplemen \\ \hline
\end{tabular}
\end{center}
\end{frame}

% ============================================================
% SECTION 9-3: Operasi dan Evaluasi
% ============================================================
\section{Operasi Dasar dan Evaluasi}

\begin{frame}
\frametitle{Tiga Operasi Fundamental}
\begin{center}
\small
\begin{tabular}{|c|c||c|c|c|}
\hline
$x$ & $y$ & $x+y$ (OR) & $x \cdot y$ (AND) & $x'$ (NOT) \\ \hline
0 & 0 & 0 & 0 & 1 \\ \hline
0 & 1 & 1 & 0 & 1 \\ \hline
1 & 0 & 1 & 0 & 0 \\ \hline
1 & 1 & 1 & 1 & 0 \\ \hline
\end{tabular}
\end{center}
\begin{block}{Prioritas}
NOT ($'$) $>$ AND ($\cdot$) $>$ OR ($+$). Contoh: $x + yz' = x + (y \cdot (z'))$
\end{block}
\end{frame}

\begin{frame}
\frametitle{Evaluasi Fungsi Boolean}
$F(x,y,z) = x \cdot (y + z')$ untuk $(1,0,1)$:\\
$F = 1 \cdot (0 + 1') = 1 \cdot (0+0) = 1 \cdot 0 = 0$\\
Fungsi $n$ variabel: $2^n$ kombinasi input.
\end{frame}

% ============================================================
% SECTION 9-4: Hukum Fundamental
% ============================================================
\section{Hukum-Hukum Fundamental}

\begin{frame}
\frametitle{Rangkuman Hukum Boolean}
\begin{center}
\small
\begin{tabular}{|l|l|l|}
\hline
\textbf{Hukum} & \textbf{OR} & \textbf{AND} \\ \hline
Identitas & $x+0=x$ & $x \cdot 1 = x$ \\ \hline
Dominasi & $x+1=1$ & $x \cdot 0 = 0$ \\ \hline
Idempoten & $x+x=x$ & $x \cdot x = x$ \\ \hline
Komplemen & $x+x'=1$ & $x \cdot x' = 0$ \\ \hline
De Morgan & $(x+y)' = x'y'$ & $(xy)' = x' + y'$ \\ \hline
Absorpsi & $x + xy = x$ & $x(x+y) = x$ \\ \hline
Konsensus & $xy + x'z + yz = xy + x'z$ & \\ \hline
\end{tabular}
\end{center}
\end{frame}

\begin{frame}
\frametitle{Prinsip Dualitas}
\begin{definisi}[Dualitas]
Setiap identitas tetap valid jika $+ \leftrightarrow \cdot$ dan $0 \leftrightarrow 1$ ditukar secara simultan.
\end{definisi}
Contoh: dual dari $x+0=x$ adalah $x \cdot 1 = x$.
\end{frame}

\begin{frame}
\frametitle{Contoh Penyederhanaan}
$(x+y)(x+y') = x + (y \cdot y') = x + 0 = x$\\
$xy + xy' = x(y+y') = x \cdot 1 = x$
\end{frame}

% ============================================================
% SECTION 9-5: Penyederhanaan
% ============================================================
\section{Penyederhanaan}

\begin{frame}
\frametitle{Strategi Penyederhanaan}
\begin{enumerate}
  \item Faktor persekutuan (distributif)
  \item Hukum komplemen untuk eliminasi
  \item Hukum absorpsi untuk menghilangkan redundan
  \item Hukum konsensus
\end{enumerate}
Tujuan: minimal literal dan operator $\Rightarrow$ sedikit gerbang.
\end{frame}

\begin{frame}
\frametitle{Contoh: $xy + xy' + x'yz$}
$F = x(y+y') + x'yz = x \cdot 1 + x'yz = x + x'yz = x + yz$ (absorpsi)
\end{frame}

% ============================================================
% SECTION 9-6: SOP dan POS
% ============================================================
\section{Bentuk Kanonik SOP dan POS}

\begin{frame}
\frametitle{Literal, Minterm, Maxterm}
\begin{block}{Literal}
Variabel atau komplemennya: $x$, $y'$, $z$
\end{block}
\begin{block}{Minterm $m_j$}
Suku perkalian dengan semua $n$ variabel (asli/komplemen). Bernilai 1 untuk tepat satu kombinasi. Contoh: $m_5 = xy'z$ (101)
\end{block}
\begin{block}{Maxterm $M_j$}
Suku penjumlahan dengan semua $n$ variabel. Bernilai 0 untuk tepat satu kombinasi. $m_j = (M_j)'$
\end{block}
\end{frame}

\begin{frame}
\frametitle{SOP dan POS}
\begin{definisi}[SOP]
$F = \sum m(j_1, j_2, \ldots)$ = OR dari minterm di mana $F=1$
\end{definisi}
\begin{definisi}[POS]
$F = \prod M(j_1, j_2, \ldots)$ = AND dari maxterm di mana $F=0$
\end{definisi}
Konversi: $F = \sum m(1,4,7)$ $\Leftrightarrow$ $F = \prod M(0,2,3,5,6)$ (indeks komplementer)
\end{frame}

% ============================================================
% SECTION 9-7: Rangkaian Digital
% ============================================================
\section{Rangkaian Digital}

\begin{frame}
\frametitle{Gerbang Logika}
\begin{center}
\small
\begin{tabular}{|l|l|l|}
\hline
\textbf{Gerbang} & \textbf{Ekspresi} & \textbf{Fungsi} \\ \hline
AND & $x \cdot y$ & Output 1 jika semua input 1 \\ \hline
OR & $x + y$ & Output 1 jika salah satu 1 \\ \hline
NOT & $x'$ & Membalik nilai \\ \hline
NAND & $(xy)'$ & NOT-AND \\ \hline
NOR & $(x+y)'$ & NOT-OR \\ \hline
XOR & $x \oplus y$ & 1 jika jumlah input-1 ganjil \\ \hline
\end{tabular}
\end{center}
\end{frame}

\begin{frame}
\frametitle{Kelengkapan Fungsional}
\begin{block}{Functionally Complete}
Himpunan $\{AND, OR, NOT\}$ lengkap. \textbf{NAND} dan \textbf{NOR} masing-masing lengkap secara individual (Sheffer 1913).
\end{block}
NOT: $x' = (xx)'$; AND: $xy = ((xy)')'$; OR: $x+y = (x'y')'$
\end{frame}

\begin{frame}
\frametitle{Contoh: Voting Mayoritas 3-Input}
$F=1$ jika $\geq 2$ dari 3 pemilih setuju.\\
$F = \sum m(3,5,6,7) = A'BC + AB'C + ABC' + ABC$\\
Sederhanakan: $F = AB + AC + BC$ (3 AND, 1 OR)
\end{frame}

% ============================================================
% RANGKUMAN
% ============================================================
\section{Rangkuman}

\begin{frame}
\frametitle{Rangkuman}
\begin{itemize}
  \item Aljabar Boolean: $\{0,1\}$, operator $+$, $\cdot$, $'$. Isomorfik dengan logika dan himpunan.
  \item Postulat Huntington: 6 aksioma dasar. Hukum lain = teorema turunan.
  \item Prinsip dualitas: tukar $+ \leftrightarrow \cdot$, $0 \leftrightarrow 1$.
  \item SOP ($\sum m$) dan POS ($\prod M$); saling konversi.
  \item NAND dan NOR masing-masing functionally complete.
\end{itemize}
\end{frame}

\begin{frame}
\frametitle{Kata Kunci}
\begin{center}
\small
aljabar Boolean $\bullet$ Boole $\bullet$ Shannon $\bullet$ Postulat Huntington $\bullet$ dualitas\\
literal $\bullet$ minterm $\bullet$ maxterm $\bullet$ SOP $\bullet$ POS $\bullet$ gerbang logika $\bullet$ NAND $\bullet$ NOR
\end{center}
\end{frame}

\end{document}
